\documentclass[10pt,a4paper]{amsart}
\usepackage{amsmath,amssymb,amsthm,graphicx}
% Retained preamble for the Stacks Project chapter divisors.tex (Divisors).
% The chapter file pinned-original-source.tex carries only \input{preamble} on its first line, so no
% preamble text lives inside the chapter itself.  The Stacks Project uses one single shared preamble for
% every chapter; the snapshot stages it verbatim in sources/stacks-chow/preamble.tex.  The block below is
% copied mechanically, line for line, from that file's lines 14-18 (the Xy-pic setup) and lines 154-217
% (the native theorem declarations together with the mathematical macros).  Nothing was added, removed or
% redefined.  The wrapper class and the amsmath/amssymb/amsthm/graphicx base packages are supplied by the
% assembler, so the class/\externaldocument part of the shared preamble is not repeated here.
% For commutative diagrams we use Xy-pic
\usepackage[all]{xy}

% We use 2cell for 2-commutative diagrams.
\xyoption{2cell}
% Theorem environments.
%
\theoremstyle{plain}
\newtheorem{theorem}[subsection]{Theorem}
\newtheorem{proposition}[subsection]{Proposition}
\newtheorem{lemma}[subsection]{Lemma}

\theoremstyle{definition}
\newtheorem{definition}[subsection]{Definition}
\newtheorem{example}[subsection]{Example}
\newtheorem{exercise}[subsection]{Exercise}
\newtheorem{situation}[subsection]{Situation}

\theoremstyle{remark}
\newtheorem{remark}[subsection]{Remark}
\newtheorem{remarks}[subsection]{Remarks}

\numberwithin{equation}{subsection}

% Macros
%
\def\lim{\mathop{\mathrm{lim}}\nolimits}
\def\colim{\mathop{\mathrm{colim}}\nolimits}
\def\Spec{\mathop{\mathrm{Spec}}}
\def\Hom{\mathop{\mathrm{Hom}}\nolimits}
\def\Ext{\mathop{\mathrm{Ext}}\nolimits}
\def\SheafHom{\mathop{\mathcal{H}\!\mathit{om}}\nolimits}
\def\SheafExt{\mathop{\mathcal{E}\!\mathit{xt}}\nolimits}
\def\Sch{\mathit{Sch}}
\def\Mor{\mathop{\mathrm{Mor}}\nolimits}
\def\Ob{\mathop{\mathrm{Ob}}\nolimits}
\def\Sh{\mathop{\mathit{Sh}}\nolimits}
\def\NL{\mathop{N\!L}\nolimits}
\def\CH{\mathop{\mathrm{CH}}\nolimits}
\def\proetale{{pro\text{-}\acute{e}tale}}
\def\etale{{\acute{e}tale}}
\def\QCoh{\mathit{QCoh}}
\def\Ker{\mathop{\mathrm{Ker}}}
\def\Im{\mathop{\mathrm{Im}}}
\def\Coker{\mathop{\mathrm{Coker}}}
\def\Coim{\mathop{\mathrm{Coim}}}

% Boxtimes
%
\DeclareMathSymbol{\boxtimes}{\mathbin}{AMSa}{"02}

%
% Macros for moduli stacks/spaces
%
\def\QCohstack{\mathcal{QC}\!\mathit{oh}}
\def\Cohstack{\mathcal{C}\!\mathit{oh}}
\def\Spacesstack{\mathcal{S}\!\mathit{paces}}
\def\Quotfunctor{\mathrm{Quot}}
\def\Hilbfunctor{\mathrm{Hilb}}
\def\Curvesstack{\mathcal{C}\!\mathit{urves}}
\def\Polarizedstack{\mathcal{P}\!\mathit{olarized}}
\def\Complexesstack{\mathcal{C}\!\mathit{omplexes}}
% \Pic is the operator that assigns to X its picard group, usage \Pic(X)
% \Picardstack_{X/B} denotes the Picard stack of X over B
% \Picardfunctor_{X/B} denotes the Picard functor of X over B
\def\Pic{\mathop{\mathrm{Pic}}\nolimits}
\def\Picardstack{\mathcal{P}\!\mathit{ic}}
\def\Picardfunctor{\mathrm{Pic}}
\def\Deformationcategory{\mathcal{D}\!\mathit{ef}}
\usepackage{cleveref}
\crefname{originaldefinition}{original definition}{original definitions}
\Crefname{originaldefinition}{Original definition}{Original definitions}
\makeatletter
\newlabel{morphisms-definition-scheme-theoretically-dense}{{01RB}{source}{}{sourcepointer-1}{}}
\newlabel{morphisms-definition-scheme-theoretically-dense@cref}{{[originaldefinition][0][]01RB}{[1][0][]source}{}{}{}}
\makeatother
\crefname{originalreference}{original reference}{original references}
\Crefname{originalreference}{Original reference}{Original references}
\makeatletter
\newlabel{morphisms-example-scheme-theoretic-closure}{{056C}{source}{}{sourcepointer-2}{}}
\newlabel{morphisms-example-scheme-theoretic-closure@cref}{{[originalreference][0][]056C}{[1][0][]source}{}{}{}}
\makeatother
\crefname{originallemma}{original lemma}{original lemmas}
\Crefname{originallemma}{Original lemma}{Original lemmas}
\makeatletter
\newlabel{properties-lemma-immersion-into-noetherian}{{01OX}{source}{}{sourcepointer-3}{}}
\newlabel{properties-lemma-immersion-into-noetherian@cref}{{[originallemma][0][]01OX}{[1][0][]source}{}{}{}}
\makeatother
\crefname{originallemma}{original lemma}{original lemmas}
\Crefname{originallemma}{Original lemma}{Original lemmas}
\makeatletter
\newlabel{algebra-lemma-qc-open}{{00F6}{source}{}{sourcepointer-4}{}}
\newlabel{algebra-lemma-qc-open@cref}{{[originallemma][0][]00F6}{[1][0][]source}{}{}{}}
\makeatother
\crefname{originalproposition}{original proposition}{original propositions}
\Crefname{originalproposition}{Original proposition}{Original propositions}
\makeatletter
\newlabel{algebra-proposition-minimal-primes-associated-primes}{{02CE}{source}{}{sourcepointer-5}{}}
\newlabel{algebra-proposition-minimal-primes-associated-primes@cref}{{[originalproposition][0][]02CE}{[1][0][]source}{}{}{}}
\makeatother
\crefname{originallemma}{original lemma}{original lemmas}
\Crefname{originallemma}{Original lemma}{Original lemmas}
\makeatletter
\newlabel{algebra-lemma-zero-at-ass-zero}{{0311}{source}{}{sourcepointer-6}{}}
\newlabel{algebra-lemma-zero-at-ass-zero@cref}{{[originallemma][0][]0311}{[1][0][]source}{}{}{}}
\makeatother
\title{An open subscheme of a locally Noetherian scheme is scheme theoretically dense if and only if it contains all embedded points, if and only if it contains all associated points}
\author{The Stacks Project Authors}
\date{stacks-project:divisors.tex:local-snapshot-20261012; GNU FDL 1.2 or later}
\begin{document}
\maketitle
\noindent\textit{Editorial note.} All delivered mathematical text is the author's own, copied from the pinned Stacks Project chapter file divisors.tex; no statement or proof step was rewritten, reordered or removed.\par\medskip
\noindent\textit{Editorial note.} The canonical Stacks tag of the counted statement is 083P (divisors-lemma-scheme-theoretically-dense-contain-embedded-points). Following the multi-claim convention for this source, source\_url is the source-level URL https://stacks.math.columbia.edu/ and the tag is carried in claim\_key and here rather than in a per-claim tag URL.\par\medskip
\noindent\textit{Editorial note.} The definition of embedded point is attached as a context span because the statement speaks of the embedded points of X without repeating the definition. The paragraph that precedes that definition in the source cites EGA and is therefore not included.\par\medskip
\noindent\textit{Editorial note.} The preamble is the chapter's shared Stacks Project preamble, retained as sources/stacks-divisors/retained-preamble.tex: a line-for-line excerpt of sources/stacks-chow/preamble.tex lines 14-18 and 154-217. No macro was invented or redefined.\par\medskip
\noindent\textit{Allowed standard prerequisites (editorial).}\begin{itemize}
\item Scheme theoretically dense open subscheme. Morphisms, Definition labelled morphisms-definition-scheme-theoretically-dense (canonical Stacks tag 01RB); cited in the statement at line 446. Supplies the notion that condition (1) of the statement refers to; the proof of the equivalence is carried out inside the counted span.
\item A quasi-compact open of a Noetherian affine scheme is a finite union of principal opens D(f\_1) up to D(f\_n). Algebra, Lemma labelled algebra-lemma-qc-open (canonical Stacks tag 00F6); invoked at proof line 457. Produces the elements f\_1, ..., f\_n whose localisations replace U throughout the proof.
\item A map A -> A\_\{f\_1\} x ... x A\_\{f\_n\} is injective exactly when every associated prime of A misses at least one f\_i. Morphisms, Example labelled morphisms-example-scheme-theoretic-closure (canonical Stacks tag 056C), used at proof line 460; and Algebra, Lemma labelled algebra-lemma-zero-at-ass-zero (canonical Stacks tag 0311), used at proof line 483. Translates scheme theoretic density into an injectivity statement and back; the two directions of the equivalence are then proved in the counted span.
\item The associated primes of a Noetherian ring are the minimal primes together with the embedded associated primes. Algebra, Proposition labelled algebra-proposition-minimal-primes-associated-primes (canonical Stacks tag 02CE); invoked at proof lines 472-475. Identifies condition (2) with condition (3) once density is read as containing all minimal primes.
\end{itemize}\medskip
\ifdefined\hypertarget\hypertarget{sourcepointer-1}{}\fi
\noindent\textit{Source reference.} Original definition 01RB: sources/stacks-chow/tags.txt, source lines 2296--2296. Cross-chapter pointer for the notion appearing in condition (1) of the statement (cited at statement line 446); the equivalence itself is proved inside the counted spans.\par\medskip
\ifdefined\hypertarget\hypertarget{sourcepointer-2}{}\fi
\noindent\textit{Source reference.} Original reference 056C: sources/stacks-chow/tags.txt, source lines 6705--6705. Cross-chapter pointer invoked at proof line 460 for the translation of scheme theoretic density into injectivity of A -> A\_\{f\_1\} x ... x A\_\{f\_n\}; both directions of that translation are then used explicitly in the counted proof.\par\medskip
\ifdefined\hypertarget\hypertarget{sourcepointer-3}{}\fi
\noindent\textit{Source reference.} Original lemma 01OX: sources/stacks-chow/tags.txt, source lines 2210--2210. Cross-chapter pointer invoked at proof line 455 for quasi-compactness of an open immersion into a Noetherian scheme; finiteness of the cover D(f\_1) union ... union D(f\_n) is then produced and used inside the counted proof.\par\medskip
\ifdefined\hypertarget\hypertarget{sourcepointer-4}{}\fi
\noindent\textit{Source reference.} Original lemma 00F6: sources/stacks-chow/tags.txt, source lines 560--560. Cross-chapter pointer invoked at proof line 457 for writing a quasi-compact open of Spec(A) as a finite union of principal opens; the resulting elements f\_1, ..., f\_n are used throughout the counted proof.\par\medskip
\ifdefined\hypertarget\hypertarget{sourcepointer-5}{}\fi
\noindent\textit{Source reference.} Original proposition 02CE: sources/stacks-chow/tags.txt, source lines 3081--3081. Cross-chapter pointer invoked at proof lines 472-475 for the description of the associated primes as the minimal primes together with the embedded associated primes; the equivalence of conditions (2) and (3) is derived from it inside the counted span.\par\medskip
\ifdefined\hypertarget\hypertarget{sourcepointer-6}{}\fi
\noindent\textit{Source reference.} Original lemma 0311: sources/stacks-chow/tags.txt, source lines 3968--3968. Cross-chapter pointer invoked at proof line 483 for injectivity of A into the product of its localisations at the associated primes; this is the last step of the counted proof and is used exactly as stated.\par\medskip
\begin{lemma}
\label{lemma-scheme-theoretically-dense-contain-embedded-points}
Let $X$ be a locally Noetherian scheme. Let $U \subset X$ be an
open subscheme. The following are equivalent
\begin{enumerate}
\item $U$ is scheme theoretically dense in $X$
(Morphisms, Definition \ref{morphisms-definition-scheme-theoretically-dense}),
\item $U$ is dense in $X$ and $U$ contains all embedded points of $X$, and
\item $U$ contains all associated points of $X$.
\end{enumerate}
\end{lemma}

\begin{proof}
The question is local on $X$, hence we may assume that $X = \Spec(A)$
where $A$ is a Noetherian ring. Then $U$ is quasi-compact
(Properties, Lemma \ref{properties-lemma-immersion-into-noetherian})
hence $U = D(f_1) \cup \ldots \cup D(f_n)$
(Algebra, Lemma \ref{algebra-lemma-qc-open}).
In this situation $U$ is scheme theoretically dense in $X$ if and only if
$A \to A_{f_1} \times \ldots \times A_{f_n}$ is injective, see
Morphisms, Example \ref{morphisms-example-scheme-theoretic-closure}.
Condition (3) translated into algebra means that for every associated
prime $\mathfrak p$ of $A$ there exists an $i$ with $f_i \not \in \mathfrak p$.

\medskip\noindent
Assume (1), i.e., $A \to A_{f_1} \times \ldots \times A_{f_n}$ is injective.
Say $a \in A$ has annihilator a prime $\mathfrak p$. Since $a$ maps
to a nonzero element of $A_{f_i}$ for some $i$ we see that
$f_i \not \in \mathfrak p$. Thus (3) holds.

\medskip\noindent
Note that $U$ is dense in $\Spec(A)$ if and only if it contains all
minimal prime ideals. Since the set of associated primes
of $A$ is equal to the union of the set of minimal primes and the
set of embedded associated primes (see Algebra, Proposition
\ref{algebra-proposition-minimal-primes-associated-primes}),
we see that (2) and (3) are equivalent.

\medskip\noindent
Assume (3). This means that every associated prime $\mathfrak p$ of $A$
corresponds to a prime of $A_{f_i}$ for some $i$. Then
$A \to A_{f_1} \times \ldots \times A_{f_n}$ is injective because
$A \to \prod_{\mathfrak p \in \text{Ass}(A)} A_\mathfrak p$ is injective
by Algebra, Lemma \ref{algebra-lemma-zero-at-ass-zero}. In this way
we see that (3) implies (1).
\end{proof}

\begin{definition}
\label{definition-embedded}
Let $X$ be a scheme.
Let $\mathcal{F}$ be a quasi-coherent sheaf on $X$.
\begin{enumerate}
\item An {\it embedded associated point} of $\mathcal{F}$
is an associated point which is not maximal among the
associated points of $\mathcal{F}$, i.e., it is the specialization
of another associated point of $\mathcal{F}$.
\item A point $x$ of $X$ is called an {\it embedded point}
if $x$ is an embedded associated point of $\mathcal{O}_X$.
\item An {\it embedded component} of $X$ is an irreducible
closed subset $Z = \overline{\{x\}}$ where $x$ is an embedded
point of $X$.
\end{enumerate}
\end{definition}

\end{document}
